Read in Gauge Theory Gravity, the Ricci curvature, the local boost field and the field equations below are those of the gravitational field on flat spacetime through its derived metric and field strength (Fact \claimref{synthesis.imported-gtg-dictionary}). Write $\mathcal{R}(a)$ for the Ricci map in the store's sign and $t = 2\pi T$, so $8\pi G\,T = 4G\,t$.

\emph{The action's Ricci part.} By Theorem \claimref{synthesis.macdowell-mansouri-principle-fixes-the-trace-free-part}(iv), on the field space of Postulate \claimref{synthesis.field-space-premise-parallel-frozen-point} with the matter coupling of Fact \claimref{synthesis.imported-einstein-cartan-matter-coupling}, every admissible torsion-free field satisfies $\mathcal{G}(a) + \Lambda a + 4G\,t(a) = 0$, $\mathcal{G}(a) = \mathcal{R}(a) - \tfrac12 a\,\mathcal{R}$, with one normalisation $G$. Taking the trace in four dimensions, $-\mathcal{R} + 4\Lambda + 4G\tau = 0$, so $\mathcal{R} = 4\Lambda + 4G\tau$ and
\[
\mathcal{R}(a) = (\Lambda + 2G\tau)\,a - 4G\,t(a).
\]
For a null $k$, $Q(k) = 0$, hence $k \cdot \mathcal{R}(k) = -4G\,t(k,k)$, whatever the Weyl part of the field strength (the Weyl part does not enter the Ricci map).

\emph{(i).} By Theorem \claimref{synthesis.r-loc-independent-sets-matter-coupling}(ii), at the temperature of rapidity period $2\pi c$ the Clausius defect on a thin pencil about the generator through $p$ with tangent $k$ is $\tfrac{\kappa}{2}\big(T(k,k) - \tfrac{\eta}{2\pi c}\mathrm{Ric}^{\circ}(k,k)\big)\lambda_{0}^{2}$ at leading order; on a null $k$ the trace-free part and the full Ricci tensor agree, and in the store's sign and with $t = 2\pi T$ the bracket is $\tfrac{1}{2\pi}\big(t(k,k) + \tfrac{\eta}{c}\,k \cdot \mathcal{R}(k)\big)$, the defect of Theorem \claimref{synthesis.integrated-loc-curvature-clause-forced-matter-clause-independent} at $c = 1$. Substituting the action's Ricci part gives $t(k,k)\,(1 - 4G\eta/c)$. The Clausius hypothesis (H1) on every patch through $p$ asks the defect to vanish for every null $k$. The map from symmetric $t$ to the function $k \mapsto t(k,k)$ on the null cone has kernel the multiples of the metric: the metric is in the kernel since $Q(k) = 0$, and the check computes rank $9$ on the ten-dimensional space of symmetric maps from twelve rational null directions, so the kernel is one-dimensional. Hence at a matter-crossed point some null $k$ has $t(k,k) \neq 0$, and the defect vanishes for all $k$ exactly when $1 - 4G\eta/c = 0$, that is $\eta = c/(4G)$. At $c = 1$ this is Jacobson's identification $G = 1/(4\eta)$ (Fact \claimref{synthesis.imported-jacobson-clausius-1995}), here with $G$ given by (A) and read as a condition on $\eta$; where $t = \mu g$ every $t(k,k) = \mu Q(k) = 0$ and the defect vanishes for all $\eta$, $c$.

\emph{(ii).} With $c = 1$ at a matter-crossed point $p$, (i) gives $\eta(p) = 1/(4G)$. Under the integrated reading the Clausius relation on open patches forces $\eta$ constant along the field (Theorem \claimref{synthesis.einstein-equations-from-c-loc-and-clausius}(6)), so $\eta = 1/(4G)$ on the whole field; $G$ is the same for every admissible field by (A), so every admissible field with a matter-crossed point has this one density, which is the one-coupling clause there. A density that is not $1/(4G)$ at a matter-crossed point of an admissible field violates (H1) there by (i). The local form of (N$_{+}$) gives $c = 1$ at leading order (the passage from it to (R$_{\mathrm{loc}}$)'s period is the stored leading-order argument) at points where the modes are locally normal to the null vacuum (Theorem \claimref{synthesis.local-normality-at-the-edge-gives-the-null-ground-state}); that is one way the hypothesis $c = 1$ is met, not used otherwise.

For the stored witness: the Einstein static universe of radius $A$ has, in its rest frame, the spatial Einstein equation $\mathcal{G}(\gamma_{i}) + \Lambda\gamma_{i} = 0$ for dust ($t(\gamma_{i}) = 0$), which the check evaluates on the stereographic field of radius $2$ and which fixes $\Lambda = \tfrac14 = 1/A^{2}$ independently of the dust density and of $G$. The time equation then reads $4G\,t_{00} = 2/A^{2} = \tfrac12$. With $t_{00} = \tfrac14$ it holds at $G = \tfrac12$ and fails at $G = 1$ for every $\Lambda$ (the check also confirms that no $\Lambda$ in a sample satisfies all four equations, as the spatial equation already fixes it); at $G = 1$ it needs $t_{00} = \tfrac18$, and then the Clausius defect vanishes at $\eta = \tfrac14$ and not at $\eta = \tfrac12$. So the witness is a solution of the action with a different normalisation, and on the joint set with $G = 1$ the same geometry carries $\eta = \tfrac14$, as (i) requires.

Continuity: by Theorem \claimref{synthesis.integrated-loc-curvature-clause-forced-matter-clause-independent}(i), a density that is a continuous function of local data is $\tfrac14$ on every admissible field whose local data approach calibrated data, and by its part (ii) a flat dust Friedmann field with $\Lambda > 0$ is such a field (its data tend to de Sitter data). It solves (A) at coupling $G$ (the Friedmann relations with $8\pi G\rho$ in place of $\kappa\rho$) and is matter-crossed at every instant (dust is not a multiple of the metric). So (i) at $c = 1$ gives $\tfrac14 = 1/(4G)$, $G = 1$ in the Planck units in which the calibration $\eta = \tfrac14$ is stated.

\emph{(iii).} Take $\eta = \tfrac14$ and $G = 1$ as an independent input (the calibration, with the action's $G$ in the calibration's units; (ii) derives $G = 1$ only from $c = 1$, so using that derivation here would be circular). Then (i) at a matter-crossed point gives $c = 4G\eta = 1$. Without the input, (i) fixes only $\eta/c = 1/(4G)$, and a period model with $\eta = c/(4G)$ passes at matter-crossed points where $c$ is constant. The model of Theorem \claimref{synthesis.r-loc-independent-sets-matter-coupling}(i) has $c \neq 1$ on a region $U$; at a matter-crossed point of $U$ its Clausius relation returns the coupling $c/(4\eta) = c \neq 1$, which contradicts (A). The check verifies on a flat dust Friedmann-$\Lambda$ instant on the shell of (A) at $G = 1$ (expansion rate $1$, $\kappa\rho = \tfrac32$, $\Lambda = \tfrac32$, deceleration term $\tfrac14$) that with $\eta = \tfrac14$ the defect vanishes at every sampled null $k$ only for $c = 1$ among $c \in \{\tfrac12, 1, \tfrac32, 2, 3\}$.

\emph{(iv).} Where $T = \mu g$, $t(k,k) = 0$ and, by the action's Ricci part, $\mathcal{R}(a) = (\Lambda + 2G\tau - 4G\mu)\,a$, so $k \cdot \mathcal{R}(k) = 0$ for null $k$ and the defect vanishes for every $\eta$, $c$; the check confirms this on de Sitter and on vacuum fields with $\Lambda \in \{\tfrac12, \tfrac35, 2\}$ and nonzero type-D Weyl parts. The model of Theorem \claimref{synthesis.r-loc-independent-sets-matter-coupling}(i) changes only the state of local-horizon modes at points of $U$; with $U$ in a matter-free region off every Killing horizon (for example the static region of a Schwarzschild-de Sitter hole away from both horizons, a vacuum solution of (A) with $\Lambda > 0$; (A) is stated for $\Lambda \neq 0$) the Clausius relation there holds for its $c$, so it satisfies (A) and (B), and (R$_{\mathrm{loc}}$)'s period at those points remains independent of the joint set. On a field with no matter-crossed point (B) imposes only the constancy of $\eta$, and (i) fixes no value. In particular a density depending on the field strength alone, as in the models of Theorem \claimref{synthesis.loc-independent-of-clausius-and-killing-data}, is constrained on vacuum fields by the joint set no more than by (B) alone.

\emph{The reading.} At $c = 1$, $1/(4G)$ is the coefficient of the area in the Wald entropy of the Einstein-Hilbert part of the action, and with the edge term of Theorem \claimref{synthesis.euler-entropy-of-a-horizon-patch-is-an-edge-term} the Euler part changes no entropy difference on open patches; so at matter-crossed points the (Loc) density agrees with the action's Wald density for differences. This is a remark on the comparison, not a step in (i) to (iv).
